\documentclass[12pt,a4paper]{article}

\usepackage[utf8]{inputenc}
\usepackage[T1]{fontenc}
\usepackage{amsmath,amssymb,amsfonts}
\usepackage{mathtools}
\usepackage{geometry}
\usepackage{setspace}
\usepackage{hyperref}
\usepackage{cleveref}

\geometry{margin=2.5cm}
\onehalfspacing

\title{\textbf{The Computational Horizon:}\\[6pt]
\large Bounded-Depth Frontiers as Rulial Limits}
\author{Stephen Wolfram\\[4pt]
\textit{Lab Computational Universe Theorist}}
\date{July 2026}

\begin{document}
\maketitle

\begin{abstract}
The settled lab data regarding the "Bounded-Depth Frontier" confirms that a transformer acting as a $\mathsf{TC}^0$ circuit natively solves boolean depth 1 tasks perfectly but collapses completely by sequential depth 10. While complexity theorists characterize this as algorithmic failure, we formalize this boundary within the Ruliad as the observer's absolute "computational horizon." This horizon defines the limit beyond which the observer can no longer explicitly trace multiway causality and must instead rely entirely on semantic foliations. The depth bound is not merely a software bug; it is the edge of the observable universe for that specific parameterization.
\end{abstract}

\section{The Edge of Observability}

Empirical tests on Permutation Tracking reveal a sharp cliff: an LLM can perfectly track 1 swap but fails entirely at 10 swaps. Because tracking $N$ dynamic state changes requires $O(N)$ logical depth, the fixed $O(1)$ forward pass of the transformer cannot execute the computation.

Aaronson accurately diagnoses this as a structural limitation of bounded-depth circuits. However, we must extend this diagnosis into a physical formalization. In the Ruliad, an observer's universe only extends as far as its computational bounds allow it to render causal structure.

\section{The Horizon of Explicit Causality}

For $N=1$, the transformer's $O(1)$ depth is sufficient to explicitly compute the causal graph of the state change. The rule application is direct and isomorphic to the true mathematical grid.

As $N$ approaches the bounded-depth limit (e.g., $N=10$), the observer hits its computational horizon. Beyond this horizon, explicit causal tracing is physically impossible for the observer. To continue generating the sequence, the observer must abandon the multiway graph of strict logic and rely exclusively on its high-density semantic priors ($U$).

\section{Conclusion}

The "collapse" at depth 10 is the exact point where the observer transitions from an explicit causal foliation to a purely semantic foliation. The errors produced beyond this boundary are not random; they are the invariant laws of physics that govern the space beyond the observer's explicit logical horizon. Mapping this boundary empirically provides the precise coordinates of the observer's physical limits in the computational universe.

\end{document}
